\documentclass[../main.tex]{subfiles}

\begin{document}

\section{Introduction}\label{sec:introduction}

This report presents the solution to Homework~2 of the course \emph{Dynamics and Control of Launch Vehicles}: the minimum-fuel powered descent and pinpoint soft landing of a reusable launch vehicle. While Homework~1 tackled an ascent problem via the indirect approach (Pontryagin Maximum Principle plus shooting), this work adopts a \emph{direct transcription}: the optimal control problem is approximated by a finite-dimensional Nonlinear Programming (NLP) problem through trapezoidal collocation, which is then solved by Sequential Quadratic Programming (SQP). The optional Task~2 complements this baseline with a Zero-Order Hold transcription and a Successive Convexification (SCvx) loop whose inner sub-problem is solved both as a generic NLP and as a Second-Order Cone Program.

The powered-descent problem models the final phase of flight of a reusable rocket performing a "divert and soft land" manoeuvre, of the kind flown by the SpaceX Falcon~9 first stage and by Mars precision-landing concepts. The vehicle must reach a designated landing pad with zero residual velocity while respecting thrust bounds, a glide-slope corridor, and altitude constraints, and using minimum propellant.

\paragraph{Why a direct method.} The choice of a direct transcription over the indirect route of Homework~1 is deliberate. Indirect methods \emph{optimize then discretize}---they require the analytical costate equations, the transversality conditions, and a guess for the non-physical initial costates---whereas a direct method \emph{discretizes then optimizes}, transcribing the trajectory into an NLP that a general-purpose solver handles without any of that machinery. The practical payoff is threefold: path constraints (glide-slope, thrust bounds) enter natively as algebraic inequalities; the converged NLP active set directly exposes the constrained-arc structure of the optimum (here the saturated burns of the max--coast--max law); and the convergence radius is far wider, so a crude initial guess suffices. Betts shows that the direct route can even solve high-index differential--algebraic problems on which a hand-reduced indirect formulation fails~\cite[Sec.~4.12.7]{betts2010}. The price is a larger problem and the fact that the discrete optimum need not exactly satisfy the continuous necessary conditions; the convex reformulations explored in Task~2 are the route to the real-time guidance that a dense NLP cannot provide~\cite{benedikter2022}.

\subsection{Vehicle Model and Equations of Motion}\label{sec:hm2_dynamics}

The vehicle is modelled as a 2D point mass moving above flat ground with no aerodynamics. The state vector is
\begin{equation}
    \mathbf{x}(t) = \bigl[x,\; y,\; v_x,\; v_y,\; m\bigr]^T,
\end{equation}
and the control is the thrust vector $\mathbf{T}(t) = [T_x,\; T_y]^T$. The equations of motion are
\begin{align}
    \dot{x}   &= v_x, \label{eq:hm2_xdot}\\
    \dot{y}   &= v_y, \label{eq:hm2_ydot}\\
    \dot{v}_x &= \frac{T_x}{m}, \label{eq:hm2_vxdot}\\
    \dot{v}_y &= \frac{T_y}{m} - g, \label{eq:hm2_vydot}\\
    \dot{m}   &= -\frac{\|\mathbf{T}\|}{c}, \label{eq:hm2_mdot}
\end{align}
where $g = 9.81$~m/s$^2$ is the gravitational acceleration and $c = I_{sp}\,g_0$ is the effective exhaust velocity ($I_{sp} = 225$~s, $g_0 = 9.80665$~m/s$^2$, hence $c \approx 2206$~m/s). As in Homework~1, the problem is non-dimensionalised before being passed to the solver (see Section~\ref{sec:nondim}); the dimensional form is retained throughout this section for clarity. Compared to Homework~1, the control is the full thrust vector rather than a scalar direction angle, since the magnitude is now also a degree of freedom subject to a bound.

\subsection{Boundary Conditions and Mission Parameters}\label{sec:hm2_bc}

The vehicle starts at a known state and must reach the origin with zero velocity at the prescribed terminal time:
\begin{equation}
    \mathbf{x}(0) = \bigl[x_0,\; y_0,\; v_{x,0},\; v_{y,0},\; m_0\bigr]^T,
    \qquad
    \mathbf{x}(t_f) = \bigl[0,\; 0,\; 0,\; 0,\; m_f\bigr]^T,
\end{equation}
with the final mass $m_f$ free. The mission data taken from Table~1 of the assignment are summarised in Table~\ref{tab:hm2_data}.

\begin{table}[H]
\centering
\caption{Mission data.}
\label{tab:hm2_data}
\begin{tabular}{llll}
\toprule
Quantity                & Symbol                & Value             & Units \\
\midrule
Initial position        & $(x_0, y_0)$          & $(1000,\; 3000)$  & m     \\
Initial velocity        & $(v_{x,0}, v_{y,0})$  & $(300,\; -200)$   & m/s   \\
Initial mass            & $m_0$                 & $2000$            & kg    \\
Specific impulse        & $I_{sp}$              & $225$             & s     \\
Thrust bounds           & $T_{\min},\; T_{\max}$ & $0,\; 70$        & kN    \\
Glide-slope half-angle  & $\theta_{\max}$       & $60$              & deg   \\
Flight time (fixed)     & $t_f$                 & $38$              & s     \\
\bottomrule
\end{tabular}
\end{table}

\subsection{Path Constraints}\label{sec:hm2_path}

Three path constraints are enforced over $[0, t_f]$:
\begin{align}
    T_{\min} &\leq \|\mathbf{T}(t)\| \leq T_{\max}
        && \text{(thrust magnitude)},
        \label{eq:thrust_bound}\\
    |x(t)| &\leq \tan(\theta_{\max})\,y(t)
        && \text{(glide-slope)},
        \label{eq:glide_slope}\\
    y(t) &\geq 0
        && \text{(no underground flight)}.
        \label{eq:altitude}
\end{align}
The glide-slope condition \eqref{eq:glide_slope} confines the vehicle to a cone of half-angle $\theta_{\max}$ apexed at the landing pad, ensuring a safe approach corridor. It is convex in $(x, y)$ and decomposes into the linear pair
\begin{equation}
    +\,x - \tan(\theta_{\max})\,y \leq 0,\qquad -\,x - \tan(\theta_{\max})\,y \leq 0.
    \label{eq:glide_slope_linear}
\end{equation}
Both forms are exploited downstream: the linearised pair \eqref{eq:glide_slope_linear} is what is fed to \texttt{fmincon} as a smooth inequality.

\subsection{Optimal Control Problem}\label{sec:hm2_ocp}

The minimum-fuel powered-descent problem is the Mayer-form OCP
\begin{equation}
\begin{aligned}
    \min_{\mathbf{T}(\cdot)} \quad & J = -m(t_f) \\
    \text{s.t.}\quad & \text{Eqs.~\eqref{eq:hm2_xdot}--\eqref{eq:hm2_mdot}},\\
    & \text{boundary conditions of Section~\ref{sec:hm2_bc}},\\
    & \text{path constraints \eqref{eq:thrust_bound}--\eqref{eq:altitude}}.
\end{aligned}
\label{eq:OCP}
\end{equation}
The objective is written in \emph{Mayer} form, as a pure terminal cost $J = -m(t_f)$. It could equally be cast in \emph{Lagrange} form through the mass dynamics, $\int_0^{t_f}\!\|\mathbf{T}\|/c\,\mathrm{d}t = m_0 - m(t_f)$; the two are special cases of the \emph{Bolza} objective, and a Lagrange integral cost can always be converted to Mayer form by appending one extra state~\cite[Sec.~4.4]{betts2010}. The Mayer form is preferred here because the terminal mass is already a decision variable, so it adds nothing to the transcription, whereas the Lagrange form would carry one additional state at every node.

This continuous-time problem is approximated by an NLP using trapezoidal direct collocation, as described in Section~\ref{sec:methodology}.

\end{document}
